\clearpage
\subsection{América del Norte: TLCAN-TMEC}


En 1986, México se adhirió al \acrfull{gatt}, hecho que marcó un punto decisivo en su proceso de apertura comercial. Después, en 1989 Canadá y Estados Unidos pusieron en vigor un tratado de libre comercio (\acrshort{cusfta}) orientado a eliminar aranceles y reducir barreras comerciales de bienes y servicios. Tanto la apertura comercial de México como la experiencia previa del \acrshort{cusfta} sentaron las bases para la negociación del \acrfull{tlcan} entre México, Estados Unidos y Canadá, que entró en vigor el 1° de enero de 1994 \parencite[p.~170]{Nagy2025}. La negociación del tratado estuvo acompañada de la firma de dos acuerdos de cooperación en materia de regulación laboral\footnote{\gls{alcan}} y ambiental,\footnote{\gls{acna}} debido a la controversia política que rodeó la ratificación del \gls{tlcan} en los Estados Unidos. Sin embargo, ambos acuerdos de cooperación no contaron con fuerza vinculante ni mecanismos de cumplimiento obligatorio \parencite[p.~277, 283]{Stanford2019}. 

El acuerdo entre Canadá y Estados Unidos y la posterior integración de México representaron uno de los primeros ejemplos de acuerdos comerciales regionales que iban más allá de la simple eliminación de aranceles, al incluir una amplia gama de disposiciones regulatorias, mecanismos de protección de la propiedad intelectual y de las inversiones, así como limitaciones a las intervenciones gubernamentales \parencite[pp.~279--281]{Stanford2019}. El documento estaba integrado por 22 capítulos en los que se abordaban diferentes puntos, como el programa de reducción arancelaria, la reducción de barreras no arancelarias, la desregulación industrial, disposiciones para la protección de la inversión extranjera intrarregional, mecanismos para la solución de controversias entre empresas y gobiernos, revisión por paneles binacionales de resoluciones antidumping, entre otros \parencite[véase][]{Stanford2019}. Dos de los mecanismos más importantes del \gls{tlcan} son las normas de \textit{trato nacional} y el \textit{trato de nación más favorecida}, que se encuentran en el capítulo 1 del acuerdo. El primer principio implica que los bienes, servicios e inversiones de empresas de un país del \gls{tlcan} no deben sufrir discriminación debido a su origen nacional. El principio del trato de nación más favorecida establece que cada país dentro del acuerdo debe extender a los otros dos integrantes el trato más favorable que ofrece a países externos al \gls{tlcan} en términos de aranceles y regulación \parencite[p.~21]{Anderson2008}.

La reducción de aranceles se desarrolló de manera gradual. Los aranceles sobre el comercio bilateral entre Canadá y Estados Unidos siguieron el esquema decenal acordado en 1989 en el \acrshort{cusfta}. Por su parte, los aranceles entre México y Estados Unidos y entre México y Canadá siguieron un programa de reducción de diez años a partir de 1993; de tal manera que, para 2003, gran parte de los aranceles sobre el comercio intrarregional se habían eliminado, de acuerdo con \textcite{Stanford2019}. Sin embargo, algunos sectores considerados estratégicos no fueron afectados por el acuerdo, como la industria petrolera mexicana, los servicios culturales canadienses y el transporte aéreo y las radiotelecomunicaciones estadounidenses \parencite[p.~170, 174]{Nagy2025}. Además de la reducción arancelaria, el tratado de libre comercio incluía unas \textit{reglas de origen} que evitaban que los exportadores de fuera de la región aprovecharan el país del \gls{tlcan} con los aranceles más bajos para acceder al mercado regional. Como regla general, al menos el 50\% del costo neto de un producto debe haberse generado dentro de la región, aunque el porcentaje es mucho más alto en la industria automotriz, de autopartes y de calzado \parencite[pp.~170,~182]{Nagy2025}.

%En el caso de algunos productos agrícolas, textiles y automotrices los aranceles se eliminaron gradualmente entre 1994 y 2008 \parencite[p.~170]{Nagy2025}. % Ver seccion 7.2 sobre industrias específicas de Nagy 2025 (174())

En materia energética, el \gls{cusfta} incluía una disposición que garantizaba el acceso estadounidense a los suministros energéticos de Canadá, la cual se mantuvo en el capítulo 6 del \gls{tlcan}. Por el contrario, el tratado protegió explícitamente el derecho de México sobre sus recursos energéticos y la propiedad estatal de \acrshort{pemex} \parencite[p.~282]{Stanford2019}. 

En 2018, se anunció un nuevo acuerdo comercial entre Estados Unidos y México, y un mes después se modificó para integrar también a Canadá. El nuevo acuerdo comercial, denominado \gls{tmec}, se firmó a finales de noviembre de 2018 y se implementó a partir de julio de 2020 \parencite[p.~172]{Nagy2025}. En contraste con el \gls{tlcan}, el \gls{tmec} refuerza la regulación laboral, incluyendo referencias a los derechos de los trabajadores a la asociación, la huelga y la representación sindical democrática \parencite[p.~283]{Stanford2019}. Sin embargo, el nuevo tratado no produjo grandes cambios que profundizaran la integración económica regional, sino que impuso nuevas restricciones y refuerzó algunas características proteccionistas del antiguo acuerdo. Por ejemplo, se elevó la regla de origen en la industria automotriz de un 62.5\% a un 75\%. Además, a diferencia del \gls{tlcan}, el \gls{tmec} tiene una cláusula que estipula la expiración del acuerdo en 16 años y una revisión cada 6 años \parencite[p.~184]{Nagy2025}.% https://www.scielo.org.mx/scielo.php?pid=S1870-35502021000200289&script=sci_abstract&tlng=en

El \gls{tlcan}-\gls{tmec} han incrementado la integración comercial de América del Norte, aunque este proceso se ha desarrollado de forma desequilibrada, centrado en el comercio de México y Canadá con Estados Unidos. Muestra de ello es que las exportaciones de ambos países se dirigen casi por completo al mercado estadounidense: 92.4\% en el caso de México y 97.4\% en el de Canadá. En contraste, las exportaciones estadounidenses destinadas a sus socios regionales se distribuyen en un 39.5\% hacia México y un 60.5\% hacia Canadá \parencite[p.44]{Murillo2022}. Estos datos ilustran cómo la integración regional no se articuló de forma equilibrada entre los tres miembros, sino que se estructuró en torno a dos vínculos bilaterales con Estados Unidos, mientras que el comercio bilateral entre México y Canadá es muy bajo. 
La centralidad estadounidense también se refleja en la composición del valor agregado regional. Una muestra de ello es que en 2015 el 12\% del valor agregado contenido en las exportaciones mexicanas procedía de Estados Unidos; en contraste, la participación mexicana en las exportaciones estadounidenses fue de apenas 1\%. De manera similar, el valor agregado estadounidense representaba el 10\% de las exportaciones canadienses, frente a una participación canadiense de solo un 2\% en las exportaciones de Estados Unidos. En consecuencia, aunque el \gls{tlcan} y el \gls{tmec} han favorecido la formación de una cadena regional de valor, la participación de los tres países en ella es desigual. De acuerdo con \textcite{Murillo2022}, Estados Unidos ocupa una posición central, abasteciendo a sus dos socios regionales de bienes intermedios que, después de ser procesados o ensamblados, son reexportados al mercado estadounidense para su consumo final. México y Canadá, en cambio, presentan una forma de participación más simple, principalmente como exportadores de bienes finales dirigidos a Estados Unidos. En particular, México presenta una integración relativamente superficial en la cadena de valor regional, caracterizada por un bajo contenido de valor agregado interno y una elevada dependencia de los insumos estadounidenses.

El \gls{tlcan} y, posteriormente, el \gls{tmec} pueden ser clasificados como un \gls{fta}, de acuerdo con la clasificación de \textcite{Balassa1961}, debido a que se eliminan las barreras (arancelarias y no arancelarias) al comercio entre los países firmantes, pero cada país integrante mantiene independencia en sus políticas comerciales con países externos. En el caso del \gls{tmec}, algunas industrias están sujetas a excepciones, reservas o disposiciones especiales que limitan la aplicación de las obligaciones generales del tratado \parencite[p.~182]{Nagy2025}. Por otra parte, la profundidad del proceso de integración económica es menor que en otras regiones, como en el sudeste asiático con la \gls{asean} o en Europa con la \gls{ue} y el \gls{eee}, debido a que ni el \gls{tlcan} ni el \gls{tmec} incluyen procesos de integración monetaria, laboral, fiscal o política, ni hay intenciones de ampliar el área con la integración de nuevos miembros \parencite[pp.~272--273]{Stanford2019}. 


\begin{figure}[H]
	\centering
	\includegraphics[width=1\textwidth]{images/ComunidadesPorPais_TLCAN.png}
	\caption{Integración de las economías del \gls{tlcan} en bloques productivos.}
	\label{ComunidadesPorPais_TLCAN}
\end{figure}
%\begin{figure}[H]
%	\centering
%	\includegraphics[width=1.1\textwidth]{images/ComunidadesPorPais_TLCAN.png}
%	\caption{Integración de las economías del \gls{tlcan} en bloques productivos$^{*}$.}
%	\label{ComunidadesPorPais_TLCAN}

%	\caption*{\footnotesize $^{*}$ Las filas representan países y las columnas años. El color de cada celda corresponde al color de la comunidad a la cual pertenece la mayoría de los sectores económicos de un país en el año correspondiente. Cada comunidad se identifica con el país del vértice con la mayor fuerza de salida.}
%\end{figure}


Los resultados de la detección de comunidades reflejan un bajo nivel de integración económica, ya que, en general, las tres economías forman bloques productivos individuales durante todo el periodo de estudio (véase la figura \ref{ComunidadesPorPais_TLCAN}). Solo algunos sectores extractivos de México y Canadá se integraron a una comunidad distinta de la predominante en su país. Durante ciertos periodos, estos sectores formaron parte del bloque productivo estadounidense. En el caso de Canadá, esto podría explicarse por los acuerdos de integración energética establecidos desde el \gls{cusfta}. Específicamente, la extracción de petróleo crudo (B06) de Canadá formó parte de la comunidad estadounidense de 1998 a 2004 y, nuevamente, en 2022, junto con el mismo sector mexicano. Asimismo, la industria canadiense de metales preciosos básicos y metales no ferrosos (CB24B) formó parte de dicha comunidad en 1995, 1996, 1999 y 2000. Además, de manera esporádica, la extracción de carbón (B06) de Canadá apareció dentro de la comunidad de Japón en 1998 y 2008 (véase la figura \ref{DiagramaSankey_TLCAN} y el apéndice \ref{Apendice3}). 


\begin{figure}[H]
	\centering
	\includegraphics[width=1\textwidth]{sankey_diagrams/north_america.png}
	\caption{Evolución de la pertenencia a comunidades de los sectores económicos de México, Estados Unidos y Canadá.}
	\label{DiagramaSankey_TLCAN}
\end{figure}


La figura \ref{Subgraph_TMEC} ilustra el grado de interdependencia productiva entre Canadá, Estados Unidos y México en 2022. Destaca la asimetría en la densidad de las aristas de mayor peso asociadas a los mercados internos (color gris): estas se concentran principalmente en Estados Unidos, en contraste con Canadá y México. Además, los flujos comerciales intrarregionales más intensos, representados mediante aristas de color coral, se observan principalmente entre Canadá y Estados Unidos, así como entre Estados Unidos y México; en cambio, no se aprecia un vínculo productivo fuerte entre México y Canadá. Estas características podrían explicar por qué el algoritmo de Leiden no agrupa a las tres economías en un único bloque productivo regional, sino que las asigna a comunidades individuales.

\begin{figure}[H]
	\centering
	\includegraphics[width=1\textwidth]{subgraphs/tmec_subrgraph.png}
	\caption{Subgráfica inducida por los vértices de México, Estados Unidos y Canadá (2022).$^{*}$}
	\label{Subgraph_TMEC}
	
		\caption*{\footnotesize $^{*}$ El tamaño de cada vértice representa su fuerza de salida; su posición se obtiene mediante el algoritmo de Fruchterman--Reingold (500 iteraciones); y su color identifica el país al que pertenece, no su comunidad. Solo se presentan las aristas cuyo peso supera el percentil 95 de la distribución regional. Las aristas grises representan transacciones domésticas y las de color coral, flujos comerciales internacionales.}
\end{figure}


% C24B (Metales preciosos basicos y otros metales no ferrosos) Canadá 1995-1996, 1999-2000 En EU
% B06 Extraccion de crudo Canadá 1998-2004 y 2022 EU
% B06 Extracción de crudo México 2022 en EU

% b05 Extraccion de carbon y lignito De canada a Jqapon 1998


%Algunos autores, como \textcite{Nagy2025} y \textcite{Kondonassis2008}, coinciden en que uno de los retos más importantes del \gls{tlcan} es que fue el primer acuerdo comercial que intentó la integración entre países con niveles de desarrollo muy diferentes.


%Además, las violaciones de estas disposiciones laborales se sujetan a mecanismos de solución de controversias mediante paneles trinacionales. Otro punto importante del nuevo tratado es una disposición centrada en la industria automotriz que exige que el 30\% del valor agregado calificado en un vehículo debe ser producido por trabajadores que ganene al menos de 16 dólares por hora, lo cual es relevante para los cálculos de las reglas de origen.
% Ver seccion 7.6 sobre TMEC Nagy 2025 (181-182))

% One of the problems with NAFTA is that integration was tried among countries of such disparate levels of development. Kondonassis2008 - 17
% NAFTA is a trade agreement and after ten years remains as such, but is limited. It falls short of a full free trade agreement because it includes much exclusion that effectively deny market access. Kondonassis2008 - 21

